\section{The information in the endpoint}\label{sec:fisher}

In this section we return to the walk of Paper~A, and the notation is local
to the section. The first step is left or right with probability $\frac12$
each, and each later step repeats the one before it with probability
$p=1-q$, where $0<q<1$. Put
\[
 \rho=1-2q,\qquad u=\frac qp,\qquad x=Nq,\qquad y=Nu=\frac{Nq}p .
\]
Let $J$ be the number of switches in the $N$ steps, $K$ the number of right
steps, $X=2K-N$ the endpoint and $W=X/N$. This is Example~\ref{ex:models}(i)
with $h=q$, so $x$ is the time horizon $T$. A path with $J$ switches has
probability $\frac12p^{N-1-J}q^J$. So $J\sim\operatorname{Bin}(N-1,q)$, the
score of the path is $(J-(N-1)q)/(pq)$, and the path carries the Fisher
information $\Var J/(pq)^2=(N-1)/(pq)$ about $q$. The endpoint carries
$I_K(q)=\sum_{k=0}^N(\partial_q\Prob(K=k))^2/\Prob(K=k)$, and we call
$e_N(q)=I_K(q)\,pq/(N-1)$ the \emph{efficiency} of the endpoint. It is the
fraction of the information in the path that is left when we only see where
the walk ends.

The endpoint sees the switches only through how far the walk spreads. If
$X$ were normal with mean $0$ and variance $Np/q$, which is close to
$\Var X$ when $x$ is large, it would carry the information
$\frac12(\partial_q\log(Np/q))^2=1/(2p^2q^2)$. This is the fraction
$1/(2(N-1)pq)\approx1/(2x)$ of the information in the path. It is only a
heuristic, but Theorem~\ref{thm:fisher} confirms the leading term $1/(2x)$.
The leading term in the Gaussian approximation of Smiga, Radaelli, Binder
and Landi~\cite{smiga2023} vanishes here, since the stationary law
$(\frac12,\frac12)$ does not depend on $q$. So in that approximation all of
the endpoint information comes from the variance term that their
leading-order formula drops.

\begin{proposition}[a correlation ratio]\label{prop:eta}
For $N\ge2$ and $0<q<1$,
\[
 e_N(q)=\frac{\Var(\E[J\mid K])}{\Var J}.
\]
For $0<k<N$ we have $\E[J\mid K=k]=\frac{d}{dt}\log S_{k,N-k}(e^t)$ at
$t=\log(q/p)$, where $S_{a,b}$ is the polynomial of
Section~\ref{sec:papersAB}. For $k\in\{0,N\}$ we have $J=0$.
\end{proposition}

\begin{proof}
For each $k$, $\Prob(K=k)$ is the sum of the probabilities of the paths
with $k$ right steps. Differentiating each term gives
$\partial_q\log\Prob(K=k)=\E[(J-(N-1)q)/(pq)\mid K=k]$. Since
$\E J=(N-1)q$, this gives $I_K(q)=\Var(\E[J\mid K])/(pq)^2$, and we divide
by $(N-1)/(pq)=\Var J/(pq)^2$. By~\cite[Thm.~2.1]{paperA}, for $0<k<N$ we
have $\Prob(K=k,J=j)=\frac12p^{N-1}W(N,k,j)u^j$, where $W(N,k,j)$ is the
number of words with $k$ letters $R$, $N-k$ letters $L$ and $j$ switches.
Summing over $j$ gives $\Prob(K=k)=\frac12p^{N-1}S_{k,N-k}(u)$. Hence
$\E[J\mid K=k]=uS'_{k,N-k}(u)/S_{k,N-k}(u)$, which is the stated
derivative. The only paths with $K\in\{0,N\}$ are the 2 constant paths.
\end{proof}

So $e_N$ is the correlation ratio of $J$ on $K$, and $e_N\le1$. Since
$\E[J\mid K]$ is the orthogonal projection of $J$ onto the functions of $K$
in $L^2$, we have $\Cov(J,g(K))=\Cov(\E[J\mid K],g(K))$ for every function
$g$. Then the Cauchy--Schwarz inequality gives
\begin{equation}\label{eq:fisher-proj}
 e_N(q)\ge\frac{\Cov(J,g(K))^2}{\Var J\,\Var g(K)}
 \qquad\text{whenever }\Var g(K)>0.
\end{equation}

\begin{proposition}[few switches]\label{prop:smallx}
For $N\ge2$ and $0<q<1$,
\[
 1-(N-2)q\le p^{N-2}\le\frac{(N-1)qp^{N-2}}{1-p^{N-1}}\le e_N(q)
 \le1-\frac{(N-2)qp^{N-3}}{2p+(N-2)q}.
\]
Hence $e_N(q)=1-\frac{N-2}2q+O_N(q^2)$ as $q\to0$ with $N$ fixed. If $q\to0$
and $Nq$ stays bounded, then $e_N(q)\to1$ if and only if $Nq\to0$.
\end{proposition}

\begin{proof}
The first 2 inequalities are Bernoulli's inequality, in the forms
$p^{N-2}\ge1-(N-2)q$ and $p^{N-1}\ge1-(N-1)q$. For the third, let
$A=\mathbf 1\{K\in\{0,N\}\}$. Only the constant paths have $K\in\{0,N\}$,
so $A=\mathbf 1\{J=0\}$. Then $\Cov(J,A)=-\E J\,\E A=-(N-1)qp^{N-1}$ and
$\Var A=p^{N-1}(1-p^{N-1})$, and \eqref{eq:fisher-proj} with $g(K)=A$ gives
the third inequality.

For the upper bound, $1-e_N=\E[\Var(J\mid K)]/\Var J$ by the law of total
variance, and $\Prob(K=k)\Var(J\mid K=k)$ is the least value of
$\E[(J-c)^2\mathbf 1\{K=k\}]$ over real $c$. For $0<k<N$ there are 2 words
with 1 switch and $N-2$ words with 2 switches~\cite[Thm.~2.1]{paperA}.
So $a=\Prob(K=k,J=1)=qp^{N-2}$ and $b=\Prob(K=k,J=2)=\frac12(N-2)q^2p^{N-3}$
do not depend on $k$. For every $c$ we have
$\E[(J-c)^2\mathbf 1\{K=k\}]\ge a(1-c)^2+b(2-c)^2\ge ab/(a+b)$, and summing
over the $N-1$ interior bins gives
\[
 \E[\Var(J\mid K)]\ge\frac{(N-1)ab}{a+b}=\frac{(N-1)(N-2)q^2p^{N-2}}{2p+(N-2)q},
\]
and we divide by $\Var J=(N-1)pq$.

For fixed $N$, $1-p^{N-1}=(N-1)q\bigl(1-\frac{N-2}2q+O(q^2)\bigr)$, so both
bounds are $1-\frac{N-2}2q+O(q^2)$. If $Nq\to0$, then $e_N\ge1-(N-2)q\to1$.
Otherwise some subsequence has $Nq\to x\in(0,\infty)$ and $N\to\infty$, and
along it the upper bound tends to $1-xe^{-x}/(2+x)<1$.
\end{proof}

The reason is that 1 switch and 2 switches both leave the endpoint uniform
on the interior bins. So the endpoint cannot tell them apart, and some
information is lost as soon as 2 switches are likely.

\begin{theorem}[many switches]\label{thm:fisher}
Let $R^2_N(q)=\Cov(J,X^2)^2/(\Var J\,\Var X^2)$ be the squared correlation
of $J$ and $X^2$. For $N\ge2$ and $0<q<1$,
\[
 R^2_N(q)=\frac1{2(N-1)pq}\cdot
 \frac{\bigl[1+\rho^N-p(1-\rho^N)/x\bigr]^2}{V},
\]
where
\[
 V=1-\frac{1+8\rho+\rho^2+8\rho^{N+1}}{4px}
 +\frac{\rho(1-\rho^N)\bigl[2(1+2\rho)(2+\rho)-\rho(1-\rho^N)\bigr]}{8p^2x^2}.
\]
There are absolute constants $C$ and $N_0$ such that for $N\ge N_0$,
$0<q\le\frac12$ and $1\le x\le N^{1/4}$,
\[
 0\le e_N(q)-R^2_N(q)\le C\Bigl(\frac1{x^3}+\frac{x^3}N\Bigr).
\]
Consequently, in the same range and with an absolute implied constant,
\[
 e_N(q)=\frac1{2x}+\frac1{4x^2}+O\Bigl(\frac1{x^3}+\frac{x^3}N\Bigr).
\]
\end{theorem}

The lower bound $e_N\ge R^2_N$ is \eqref{eq:fisher-proj} with $g(K)=X^2$,
and it holds for all $N\ge2$ and $0<q<1$. The closed form and the upper
bound are proved in Appendix~\ref{app:fisher}. The quadratic does well
because Proposition~\ref{prop:eta} and Paper~A's uniform Bessel estimate for
every bin~\cite[Thm.~6.2]{paperA} make $\E[J\mid K]$ close to
$y\sqrt{1-W^2}+\frac12\approx y(1-W^2/2)+\frac12$. We do not compute $C$,
and the error term we prove is small only for very large $N$, but
numerically $e_N-R^2_N$ is already of order $x^{-3}$ at $N=50$
(Table~\ref{tab:fisher}).

\begin{proof}[Proof of the expansion in Theorem~\ref{thm:fisher}]
If $x\le18$, then $\lvert e_N-\frac1{2x}-\frac1{4x^2}\rvert\le1\le18^3x^{-3}$,
since $e_N$ and $\frac1{2x}+\frac1{4x^2}$ lie in $[0,1]$. Let $x\ge18$. Then $p\ge\frac12$,
$0\le\rho<1$ and $\rho^N\le e^{-2x}\le x^{-2}$. The middle term of $V$ lies
in $[0,9/x]$, and since $1+8\rho+\rho^2=10-20q+4q^2$ it equals
$\frac5{2x}+O(q/x+x^{-2})$. The last term lies in $[0,9/x^2]$. So
$V\ge\frac12$, $1/V=1+\frac5{2x}+O(q/x+x^{-2})$, the squared bracket is
$1-\frac2x+O(q/x+x^{-2})$, and $1/(2(N-1)pq)=(1+O(q+1/N))/(2x)$. Multiplying,
$R^2_N=\frac1{2x}+\frac1{4x^2}+O(x^{-3}+q/x+1/(Nx))$, and every implied
constant is absolute. Since $q/x=1/N\le x^3/N$, the bound on $e_N-R^2_N$
gives the expansion.
\end{proof}

\begin{corollary}[at Kagey's crossing]\label{cor:fisher-crossing}
Let $p_N$ be the value of $p$ at which the middle bin and the endpoints are
equally likely (Section~\ref{sec:papersAB}), and put $q_N=1-p_N$ and
$L=\log N$. Then
\[
 e_N(q_N)=\frac1{2L}+\frac{\log L+1+\log(8/\pi)}{4L^2}
 +O\Bigl(\frac{(\log L)^2}{L^3}\Bigr).
\]
\end{corollary}

\begin{proof}
By~\cite[Thm.~5.1]{paperA}, $p_N\ge\frac23$, so $q_N\le\frac13$. By
\cite[Thm.~5.2, Eq.~(7)]{paperA},
$x_N=Nq_N=L-\frac12\log L+c_0+O(\log L/L)$.
So $1\le x_N\le N^{1/4}$ for large $N$, and Theorem~\ref{thm:fisher} gives
$e_N(q_N)=\frac1{2x_N}+\frac1{4x_N^2}+O(L^{-3})$. Write
$x_N=L(1-\varepsilon)$ with $\varepsilon=(\frac12\log L-c_0)/L+O(\log L/L^2)$.
Then $\frac1{2x_N}=\frac1{2L}+\frac{\log L-2c_0}{4L^2}+O((\log L)^2/L^3)$ and
$\frac1{4x_N^2}=\frac1{4L^2}+O(\log L/L^3)$. Since $-2c_0=\log(8/\pi)$,
adding the 2 gives the claim.
\end{proof}

So $e_N\sim1/(2\log N)$ at Kagey's crossing, but the limit is approached
slowly. The relative correction is of order $(\log\log N)/\log N$, and
Table~\ref{tab:fisher} shows that it is large at every $N$ we computed.
The product $2e_N\log N$ is $1.74$ at $N=100$ and still $1.35$ at
$N=3200$. Both $R^2_N$ and
$\frac1{2x_N}+\frac1{4x_N^2}$ lie below $e_N$ by an amount comparable to
$x_N^{-3}$.

\begin{table}[ht]
\centering\small
\begin{tabular}{@{}rcccccc@{}}
\toprule
$N$ & $x_N$ & $e_N(q_N)$ & $2x_Ne_N$ & $2Le_N$ & $R^2_N(q_N)$
 & $\frac1{2x_N}+\frac1{4x_N^2}$\\
\midrule
50 & 2.8998 & 0.24117 & 1.3987 & 1.8869 & 0.20873 & 0.20216\\
100 & 3.5034 & 0.18869 & 1.3221 & 1.7379 & 0.16642 & 0.16309\\
400 & 4.7501 & 0.12681 & 1.2047 & 1.5196 & 0.11715 & 0.11634\\
1600 & 6.0246 & 0.09433 & 1.1366 & 1.3919 & 0.09007 & 0.08988\\
3200 & 6.6684 & 0.08361 & 1.1151 & 1.3496 & 0.08070 & 0.08060\\
\bottomrule
\end{tabular}
\caption{The efficiency of the endpoint at Kagey's crossing $q_N=1-p_N$,
computed from the exact law, with $x_N=Nq_N$ and $L=\log N$. The last 2
columns are the squared correlation of Theorem~\ref{thm:fisher} and its
two-term expansion.}
\label{tab:fisher}
\end{table}

\begin{remark}[fixed $q$ and fixed $x$]\label{rem:fisher-regimes}
Let $Nq\to\infty$ with $q\le q_0$ for a fixed $q_0<1$. Then
$\liminf2(N-1)pq\,e_N(q)\ge1$. Indeed $e_N\ge R^2_N$, and
$2(N-1)pq\,R^2_N$ is the squared bracket over $V$. Here
$px\ge(1-q_0)x\to\infty$, and $\lvert\rho\rvert^N\to0$ because
$\rho^N\le e^{-2x}$ when $q\le\frac12$ and $\lvert\rho\rvert\le2q_0-1$ when
$q>\frac12$. The middle and last terms of $V$ are at most $18/(4px)$ and
$40/(8p^2x^2)$ in absolute value, so $V\to1$. The bracket tends to $1$
since $p(1-\rho^N)/x\le2/x$.
We have not proved the matching upper bound, but numerically the ratio
tends to $1$. At $q=0.1$ it is $1.0113$, $1.0055$ and $1.0027$ for $N=400$,
$800$ and $1600$. So at fixed $q$ the efficiency is about $1/(2(N-1)pq)$,
and $1/(2Nq)$ is too small by the factor $1/p$.

If instead $x=Nq$ is fixed and $N\to\infty$, Proposition~\ref{prop:smallx}
gives $x/(e^x-1)\le\liminf e_N\le\limsup e_N\le1-xe^{-x}/(2+x)$. We do not
prove that $e_N$ converges. Numerically it is close to the same efficiency
$e_\infty(x)$ for the telegraph process on $[0,1]$ that changes direction at
rate $x$ and is observed through the time it spends moving right. Along the
crossing, $e_N/e_\infty(x_N)$ is $1.025$ at $N=100$ and $1.002$ at
$N=3200$. The product $2xe_\infty(x)$ is $1.18$, $1.40$ and $1.34$ at $x=1$,
$2$ and $3$, so for moderate $x$ neither $1$ nor $1/(2x)$ is a good
approximation.
\end{remark}
